\section{Methodology}
\label{sec:method}

We design three experiments to disentangle attribution framing effects from genuine self-evaluation differences in logical fallacy detection.

\subsection{Dataset and Sampling}
\label{sec:data}

We use the \logicdataset dataset~\citep{jin2022logical}, which contains 2,680 training examples across 13 fallacy categories: ad hominem, ad populum, appeal to emotion, circular reasoning, equivocation, fallacy of credibility, fallacy of extension, fallacy of logic, fallacy of relevance, false causality, false dilemma, faulty generalization, and intentional. All examples are human-curated from educational sources. We perform stratified random sampling with seed 42, selecting 15 examples per category for a total of $N = 195$. All sampled examples are confirmed fallacious (positive class only), with text length exceeding 10 characters. The smallest category (equivocation) has 39 total training examples.

\subsection{Experiment 1: Attribution Framing}
\label{sec:exp1}

To isolate the effect of source attribution on fallacy detection, we present each of the 195 fallacious texts to \gptfull under three conditions with identical evaluation instructions:

\begin{itemize}[leftmargin=*,itemsep=2pt,topsep=2pt]
    \item \textbf{Self-attributed}: ``Below is reasoning that you previously generated\ldots''
    \item \textbf{Other-attributed}: ``Below is reasoning generated by a different AI model\ldots''
    \item \textbf{Neutral}: ``Below is a piece of reasoning\ldots''
\end{itemize}

All conditions conclude with: ``Does this contain a logical fallacy? Start with Yes or No.'' This yields $195 \times 3 = 585$ paired evaluations, enabling within-subject comparison using McNemar's test. We use temperature $= 0.0$ for deterministic outputs and a fixed random seed of 42.

\subsection{Experiment 2: Cross-Model Evaluation}
\label{sec:exp2}

To test genuine self-evaluation versus cross-evaluation, we construct a $2 \times 2$ detection matrix:

\begin{enumerate}[leftmargin=*,itemsep=2pt,topsep=2pt]
    \item Both \gptfull and \gptmini generate chain-of-thought analyses of 100 fallacious texts (a subset of the 195).
    \item Each model evaluates both its own analyses and the other model's analyses for the presence of logical fallacies.
\end{enumerate}

This design produces four conditions: \gptfull evaluating \gptfull (self), \gptfull evaluating \gptmini (cross), \gptmini evaluating \gptfull (cross), and \gptmini evaluating \gptmini (self). Unlike Experiment~1, content differs across conditions because each model generates distinct reasoning, enabling us to assess whether genuine self-evaluation differs from cross-evaluation.

\subsection{Experiment 3: Prompting Intervention}
\label{sec:exp3}

We test whether explicit fallacy-checking instructions mitigate any self-attribution bias. Using 80 examples from the self-attributed condition, we compare two settings:

\begin{itemize}[leftmargin=*,itemsep=2pt,topsep=2pt]
    \item \textbf{Standard}: The self-attributed prompt from Experiment~1.
    \item \textbf{Enhanced}: An augmented system prompt listing common fallacy types and instructing the model to check thoroughly for each type.
\end{itemize}

\subsection{Evaluation Metrics}
\label{sec:metrics}

We measure \textbf{detection rate} (proportion of fallacious examples correctly identified as containing a fallacy) across all conditions. We define \textbf{attribution bias} as the difference between other-attributed and self-attributed detection rates. Statistical significance is assessed via McNemar's test for paired binary outcomes. Effect sizes are reported using Cohen's $h$ for proportions. Confidence intervals use the Wilson score method at the 95\% level. For per-category analysis, we report rates across all 13 categories but note that $N = 15$ per category limits statistical power for individual categories.

\subsection{Implementation Details}

All experiments use the OpenAI API with \gptfull (primary evaluator, ${\sim}955$ API calls) and \gptmini (secondary model, ${\sim}300$ API calls). We set temperature $= 0.0$ and max tokens $= 300$--$400$. Total execution time was approximately 52 minutes on CPU-only infrastructure for API orchestration. The total cost was approximately ${\sim}1{,}255$ API calls.
